\chapter{常用函数}\label{ch:g10:reffunc}

几个简单\omterm{def:g10:functions:function}{函数}——\omterm{def:g10:reffunc:affine}{一次函数}、平方\omterm{def:g10:functions:function}{函数}、反比例\omterm{def:g10:functions:function}{函数}、平方根\omterm{def:g10:functions:function}{函数}、
立方\omterm{def:g10:functions:function}{函数}——单独出现或相互组合，随处可见。熟记它们的\omterm{def:g10:functions:graph}{图像}与
变化规律，许多问题便能迅速画出草图。
本章逐一研究这些\omterm{def:g10:functions:function}{函数}，并借助
\cref{ch:g10:functions}中的比较技巧证明其变化规律。

\section{一次函数}

\begin{definition}[一次函数]\label{def:g10:reffunc:affine}
形如下式的\omterm{def:g10:functions:function}{函数}称为\emph{一次函数}\index{一次函数}：
\[
f(x) = mx + p,
\]
其中 $m$ 与 $p$ 为固定的\omterm{def:g10:numbers:sets}{实数}。其\omterm{def:g10:functions:graph}{图像}是一条直线：
$m$ 称为\emph{斜率}\index{斜率}，$p$ 称为
\emph{纵截距}\index{纵截距}（亦称 $y$ 截距，直线与 $y$ 轴相交于
点 $(0, p)$）。当 $p = 0$ 时，\omterm{def:g10:functions:function}{函数}
为\emph{正比例\omterm{def:g10:functions:function}{函数}}：$f(x) = mx$ 表示正比例关系。
\end{definition}

\begin{proposition}[斜率与单调性]\label{prop:g10:reffunc:affine}
设 $f(x) = mx + p$。
\begin{enumerate}
  \item 对任意两个不同的输入 $u \neq v$：
        $m = \dfrac{f(v) - f(u)}{v - u}$（\omterm{def:g10:reffunc:affine}{斜率}是
        输入每改变一个单位时输出的改变量）。
  \item 若 $m > 0$，则 $f$ 在 $\R$ 上严格\omterm{def:g10:functions:variations}{递增}；若 $m < 0$，
        则严格\omterm{def:g10:functions:variations}{递减}；若 $m = 0$，则为常值\omterm{def:g10:functions:function}{函数}。
\end{enumerate}
\end{proposition}

\begin{proof}
1. 计算 $f(v) - f(u) = (mv + p) - (mu + p) = m(v - u)$，再除以
$v - u \neq 0$。

2. 若 $u < v$，则 $f(v) - f(u) = m(v-u)$ 的符号与 $m$ 相同，因为
$v - u > 0$。$m$ 为正时得 $f(u) < f(v)$（\omterm{def:g10:functions:variations}{递增}），$m$ 为负时
得 $f(u) > f(v)$（\omterm{def:g10:functions:variations}{递减}）。
\end{proof}

\begin{omfigure}
\begin{tikzpicture}
\begin{axis}[omaxis,
  width=7.8cm, height=5.6cm,
  xmin=-3.4, xmax=3.4, ymin=-2.8, ymax=4.4,
  xtick={-2,-1,1,2}, ytick={-1,1,2,3}]
  \addplot[omDef, thick, domain=-3:2.6] {0.5*x + 1};
  \addplot[omThm, thick, domain=-2.2:3] {-x + 1};
  \node[omDef, font=\small] at (axis cs:2.5,2.7) {$y = \tfrac12 x + 1$};
  \node[omThm, font=\small] at (axis cs:-2.0,2.6) {$y = -x + 1$};
  \draw[dashed] (axis cs:1,1.5) -- (axis cs:2,1.5) -- (axis cs:2,2);
  \node[font=\footnotesize, below] at (axis cs:1.5,1.5) {$1$};
  \node[font=\footnotesize, right] at (axis cs:2,1.75) {$m$};
\end{axis}
\end{tikzpicture}

{\small 两个\omterm{def:g10:reffunc:affine}{一次函数}：$m = \frac12 > 0$ 时\omterm{def:g10:functions:variations}{递增}，
$m = -1 < 0$ 时\omterm{def:g10:functions:variations}{递减}。向右移动一个单位，$y$ 改变 $m$。}
\end{omfigure}

\begin{example}\label{ex:g10:reffunc:affine}
求\omterm{def:g10:functions:graph}{图像}经过点 $A(1, 3)$ 与
$B(4, 9)$ 的\omterm{def:g10:reffunc:affine}{一次函数}。\omterm{def:g10:reffunc:affine}{斜率}为
\[
m = \frac{9 - 3}{4 - 1} = \frac{6}{3} = 2 .
\]
于是 $f(x) = 2x + p$，由 $f(1) = 3$ 得 $2 + p = 3$，故 $p = 1$：
$f(x) = 2x + 1$。验证点 $B$：$f(4) = 9$。
\end{example}

\section{平方函数}

\begin{proposition}[平方函数]\label{prop:g10:reffunc:square}
\omterm{def:g10:functions:function}{函数} $f(x) = x^2$ 定义在 $\R$ 上，具有以下性质：
\begin{enumerate}
  \item 在 $\intoc{-\infty}{0}$ 上严格\omterm{def:g10:functions:variations}{递减}，在
        $\intco{0}{+\infty}$ 上严格\omterm{def:g10:functions:variations}{递增}，\omterm{def:g10:functions:extrema}{最小值}为 $0$，在 $x = 0$ 处取得；
  \item 满足 $f(-x) = f(x)$，对一切 $x$ 成立：其\omterm{def:g10:functions:graph}{图像}即
        \emph{抛物线}\index{抛物线}，
        关于纵轴对称。
\end{enumerate}
\end{proposition}

\begin{proof}
1. 设 $0 \leq u < v$，则
\[
f(v) - f(u) = v^2 - u^2 = (v - u)(v + u) > 0,
\]
因为 $v - u > 0$ 且 $v + u > 0$，故 $f$ 在
$\intco{0}{+\infty}$ 上严格\omterm{def:g10:functions:variations}{递增}。若 $u < v \leq 0$，则 $v - u > 0$，但
$v + u < 0$，从而 $f(v) - f(u) < 0$，故严格\omterm{def:g10:functions:variations}{递减}。最后
$x^2 \geq 0 = f(0)$，对所有 $x$ 成立。

2. $(-x)^2 = x^2$，故点 $(x, x^2)$ 与 $(-x, x^2)$ 关于纵轴
对称，二者都在\omterm{def:g10:functions:graph}{图像}上。
\end{proof}

\begin{remark}[平方与大小关系]
由于平方\omterm{def:g10:functions:function}{函数}在负半轴上\omterm{def:g10:functions:variations}{递减}，取平方
\emph{颠倒}负数的大小次序：$-3 < -2$，但 $9 > 4$。
两边平方前，务必先检查符号。
\end{remark}

\section{反比例函数}

\begin{proposition}[反比例函数]\label{prop:g10:reffunc:inverse}
\omterm{def:g10:functions:function}{函数} $f(x) = \dfrac1x$ 定义在 $x \neq 0$ 上，具有以下性质：
\begin{enumerate}
  \item 在 $\intoo{-\infty}{0}$ 上严格\omterm{def:g10:functions:variations}{递减}，并在
        $\intoo{0}{+\infty}$ 上严格\omterm{def:g10:functions:variations}{递减}；
  \item 满足 $f(-x) = -f(x)$：其\omterm{def:g10:functions:graph}{图像}即
        \emph{双曲线}\index{双曲线}，关于原点对称。
\end{enumerate}
\end{proposition}

\begin{proof}
1. 设 $0 < u < v$，则
\[
f(v) - f(u) = \frac1v - \frac1u = \frac{u - v}{uv} < 0,
\]
由 $u - v < 0$ 且 $uv > 0$，知\omterm{def:g10:functions:function}{函数}在
$\intoo{0}{+\infty}$ 上严格\omterm{def:g10:functions:variations}{递减}。在 $\intoo{-\infty}{0}$ 上，同一商式的
$u - v < 0$ 与 $uv > 0$ 再次成立（两个负数之积），故 $f$
在该\omterm{def:g10:numbers:interval}{区间}上同样\omterm{def:g10:functions:variations}{递减}。

2. $\frac{1}{-x} = -\frac1x$.
\end{proof}

\begin{remark}
注意：$\frac1x$ \emph{并非}在其整个\omterm{def:g10:functions:function}{定义域}上\omterm{def:g10:functions:variations}{递减}。
取 $u = -1$，$v = 1$，则\omterm{def:g10:functions:function}{函数}值由 $-1$ 跃升至 $1$。两个分支
必须分别研究。
\end{remark}

\begin{omfigure}
\begin{tikzpicture}
\begin{axis}[omaxis,
  width=5.6cm, height=5.2cm,
  xmin=-3.2, xmax=3.2, ymin=-2.6, ymax=6.4,
  xtick={-2,-1,1,2}, ytick={1,4},
  title={$y = x^2$}]
  \addplot[omDef, thick, domain=-2.45:2.45, samples=90] {x^2};
\end{axis}
\end{tikzpicture}\hfill
\begin{tikzpicture}
\begin{axis}[omaxis,
  width=5.6cm, height=5.2cm,
  xmin=-3.4, xmax=3.4, ymin=-3.4, ymax=3.4,
  xtick={-1,1}, ytick={-1,1},
  title={$y = \frac1x$}]
  \addplot[omDef, thick, domain=0.3:3.1, samples=80] {1/x};
  \addplot[omDef, thick, domain=-3.1:-0.3, samples=80] {1/x};
\end{axis}
\end{tikzpicture}

{\small 抛物线 $y = x^2$（关于纵轴对称）与
双曲线 $y = \frac1x$（两支，关于原点对称）。}
\end{omfigure}

\section{平方根与立方}

\begin{proposition}[平方根函数]\label{prop:g10:reffunc:sqrt}
\omterm{def:g10:functions:function}{函数} $f(x) = \sqrt{x}$ 定义在 $\intco{0}{+\infty}$ 上，
严格\omterm{def:g10:functions:variations}{递增}。
\end{proposition}

\begin{proof}
设 $0 \leq u < v$。分子分母同乘以\emph{共轭因式}：
\[
\sqrt v - \sqrt u
= \frac{(\sqrt v - \sqrt u)(\sqrt v + \sqrt u)}{\sqrt v + \sqrt u}
= \frac{v - u}{\sqrt v + \sqrt u} > 0,
\]
因 $v - u > 0$ 且 $\sqrt v + \sqrt u > 0$（注意 $v > 0$，故
$\sqrt v > 0$）。
\end{proof}

\begin{proposition}[立方函数]\label{prop:g10:reffunc:cube}
\omterm{def:g10:functions:function}{函数} $f(x) = x^3$ 定义在 $\R$ 上，严格\omterm{def:g10:functions:variations}{递增}，且
其\omterm{def:g10:functions:graph}{图像}关于原点对称。
\end{proposition}
\admitted

\begin{omfigure}
\begin{tikzpicture}
\begin{axis}[omaxis,
  width=5.6cm, height=5.0cm,
  xmin=-0.7, xmax=5.4, ymin=-0.7, ymax=2.9,
  xtick={1,4}, ytick={1,2},
  title={$y = \sqrt x$}]
  \addplot[omDef, thick, domain=0:5.1, samples=100] {sqrt(x)};
\end{axis}
\end{tikzpicture}\hfill
\begin{tikzpicture}
\begin{axis}[omaxis,
  width=5.6cm, height=5.0cm,
  xmin=-2.2, xmax=2.2, ymin=-3.4, ymax=3.4,
  xtick={-1,1}, ytick={-1,1},
  title={$y = x^3$}]
  \addplot[omDef, thick, domain=-1.48:1.48, samples=90] {x^3};
\end{axis}
\end{tikzpicture}

{\small 平方根（仅对 $x \geq 0$ 有定义）与立方，两者
均\omterm{def:g10:functions:variations}{递增}。}
\end{omfigure}

\begin{proposition}[比较 $x$、$x^2$ 与 $\sqrt x$]\label{prop:g10:reffunc:compare}
当 $0 < x < 1$ 时：\ \ $x^2 < x < \sqrt x$。\quad
当 $x > 1$ 时：\ \ $\sqrt x < x < x^2$。\quad
当 $x = 0$ 与 $x = 1$ 时，三者的值相等。
\end{proposition}

\begin{proof}
设 $0 < x < 1$。将 $x < 1$ 两边同乘 $x > 0$，得 $x^2 < x$。
其次证明 $x < \sqrt x$：由于两边均为正，且在正数上平方是
\omterm{def:g10:functions:variations}{递增}的，故该不等式等价于
$x^2 < x$（已证）。当 $x > 1$ 时，将 $x > 1$ 两边同乘 $x$，
得 $x^2 > x$，同样由上述平方论证可得 $\sqrt x < x$。
\end{proof}

\begin{omfigure}
\begin{tikzpicture}
\begin{axis}[omaxis,
  width=7.4cm, height=5.8cm,
  xmin=-0.35, xmax=2.5, ymin=-0.35, ymax=2.9,
  xtick={1,2}, ytick={1,2}]
  \addplot[omThm, thick, domain=0:1.65, samples=80] {x^2};
  \addplot[omDef, thick, domain=0:2.35] {x};
  \addplot[omMeth, thick, domain=0:2.35, samples=80] {sqrt(x)};
  \fill (axis cs:1,1) circle (1.5pt);
  \node[omThm, font=\small, right] at (axis cs:1.45,2.4) {$y = x^2$};
  \node[omDef, font=\small, right] at (axis cs:2.0,1.85) {$y = x$};
  \node[omMeth, font=\small, right] at (axis cs:2.0,1.28) {$y = \sqrt x$};
\end{axis}
\end{tikzpicture}

{\small 三条曲线在 $(0,0)$ 与 $(1,1)$ 处相交，并在这两点处交换次序：
在 $0$ 与 $1$ 之间，$x^2$ 最小，$\sqrt x$ 最大；
超过 $1$ 后，次序颠倒。}
\end{omfigure}

\begin{method}[比较像的大小]\label{met:g10:reffunc:compare}
要比较 $f(a)$ 与 $f(b)$ 的大小，其中 $f$ 为参考\omterm{def:g10:functions:function}{函数}：
\begin{enumerate}
  \item 将 $a$ 与 $b$ 置于一个
        $f$ 的单调性已知的\omterm{def:g10:numbers:interval}{区间}内；
  \item 若 $f$ 在该\omterm{def:g10:numbers:interval}{区间}上\omterm{def:g10:functions:variations}{递增}，则像的次序与
        自变量的次序相同；若 $f$ \omterm{def:g10:functions:variations}{递减}，则次序相反；
  \item 若 $a$ 与 $b$ 不在同一单调\omterm{def:g10:numbers:interval}{区间}内，应先
        处理符号（例如对平方\omterm{def:g10:functions:function}{函数}，比较一个负输入与一个正
        输入）。
\end{enumerate}
\end{method}

\begin{example}
比较 $\dfrac{1}{2.3}$ 与 $\dfrac{1}{2.4}$：两个输入都在
$\intoo{0}{+\infty}$ 内，反比例\omterm{def:g10:functions:function}{函数}\omterm{def:g10:functions:variations}{递减}，且
$2.3 < 2.4$，因此 $\dfrac{1}{2.3} > \dfrac{1}{2.4}$。比较 $(-1.2)^2$
与 $(-1.5)^2$：在 $\intoc{-\infty}{0}$ 上平方\omterm{def:g10:functions:function}{函数}\omterm{def:g10:functions:variations}{递减}，且
$-1.5 < -1.2$，因此 $(-1.5)^2 > (-1.2)^2$——事实上 $2.25 > 1.44$。
\end{example}

\section{练习}

\begin{exercise}[$\star$]\label{exo:g10:reffunc:1}
对每个\omterm{def:g10:reffunc:affine}{一次函数}，求其\omterm{def:g10:reffunc:affine}{斜率}与\omterm{def:g10:reffunc:affine}{纵截距}（即\omterm{def:g10:functions:graph}{图像}与 $y$ 轴交点的纵坐标），
并说明它是\omterm{def:g10:functions:variations}{递增}还是\omterm{def:g10:functions:variations}{递减}：
\[
f(x) = 3x - 2, \qquad
g(x) = -\tfrac12 x + 4, \qquad
h(x) = 7, \qquad
k(x) = -x .
\]
\end{exercise}

\begin{exercise}[$\star$]\label{exo:g10:reffunc:2}
求\omterm{def:g10:functions:graph}{图像}经过点 $(2, 1)$ 与
$(5, 10)$ 的\omterm{def:g10:reffunc:affine}{一次函数}；再求过 $(-1, 4)$ 与 $(3, -4)$ 的\omterm{def:g10:reffunc:affine}{一次函数}。
\end{exercise}

\begin{exercise}[$\star$]\label{exo:g10:reffunc:3}
不用计算器，比较下列各组数的大小：
\[
(3.1)^2 \text{ 与 } (3.2)^2; \qquad
(-2.7)^2 \text{ 与 } (-2.8)^2; \qquad
\frac{1}{5.1} \text{ 与 } \frac{1}{5.2}; \qquad
\sqrt{17} \text{ 与 } \sqrt{15}.
\]
\end{exercise}

\begin{exercise}[$\star$]\label{exo:g10:reffunc:4}
利用平方\omterm{def:g10:functions:function}{函数}的\omterm{def:g10:functions:graph}{图像}，解\omterm{def:g10:algebra:equation}{方程} $x^2 = 16$，然后
解不等式 $x^2 \leq 16$，再解不等式 $x^2 > 9$。
\end{exercise}

\begin{exercise}[$\star$]\label{exo:g10:reffunc:5}
设 $x \in \intoo{0}{1}$。将数 $x$、$x^2$、$x^3$ 与
$\sqrt x$ 从小到大排列，并
用 $x = 0.25$ 检验你的结果。
\end{exercise}

\begin{exercise}[$\star\star$]\label{exo:g10:reffunc:6}
先用\omterm{def:g10:functions:graph}{图像}法、再用代数法求解：$\dfrac1x = 2$，以及
$\dfrac1x < 2$ 在 $x > 0$ 时的解。若再允许 $x < 0$，结果有何变化？
\end{exercise}

\begin{exercise}[$\star\star$]\label{exo:g10:reffunc:7}
已知平方\omterm{def:g10:functions:function}{函数}在 $\intoc{-\infty}{0}$ 上\omterm{def:g10:functions:variations}{递减}，
在 $\intco{0}{+\infty}$ 上\omterm{def:g10:functions:variations}{递增}，求 $x^2$ 在下列各情形下的范围：
\[
\text{(a) } x \in \intcc{2}{5};
\qquad
\text{(b) } x \in \intcc{-3}{-1};
\qquad
\text{(c) } x \in \intcc{-2}{3}.
\]
（在情形 (c) 中要小心：$x^2$ 的\omterm{def:g10:functions:extrema}{最小值}不在端点处。）
\end{exercise}

\begin{exercise}[$\star\star$]\label{exo:g10:reffunc:8}
某出租车公司收取固定费用 $4$，另加每千米 $1.5$；
另一家公司不收固定费用，每千米收 $2$。
用距离 $x$ 的\omterm{def:g10:reffunc:affine}{一次函数}分别表示两家公司的收费，作出两条直线，并求
从何距离起第一家公司更便宜。
\end{exercise}

\begin{exercise}[$\star\star$]\label{exo:g10:reffunc:9}
证明：对所有 $a, b \geq 0$，有 $\sqrt{a + b} \leq \sqrt a + \sqrt b$。
（提示：两边均非负，故可比较它们的平方。）
等号何时成立？
\end{exercise}

\begin{exercise}[$\star\star\star$]\label{exo:g10:reffunc:10}
设 $f(x) = \dfrac{2x + 1}{x - 1}$，其在 $x \neq 1$ 时有定义。
\begin{enumerate}
  \item 证明：$f(x) = 2 + \dfrac{3}{x - 1}$ 对所有 $x \neq 1$ 成立。
  \item 利用反比例\omterm{def:g10:functions:function}{函数}的单调性，推出 $f$ 在 $\intoo{1}{+\infty}$
        上的单调性。
\end{enumerate}
\end{exercise}

\section{问题：自然规律用参考函数表述}

\begin{problem}[{周末问题——刹车距离随
$v^2$ 增长，杠杆遵循 $1/d$，行星服从 $a^{3/2}$：
用参考\omterm{def:g10:functions:function}{函数}读懂物理}]
\label{pb:g10:reffunc:1}
翻开一本物理书，每一页都会出现同样的几类\omterm{def:g10:functions:graph}{图像}：
抛物线、双曲线、平方根曲线。自然界
正是用本章的参考\omterm{def:g10:functions:function}{函数}书写其规律——而熟悉
它们的形状（\cref{prop:g10:reffunc:square}、
\cref{prop:g10:reffunc:inverse}、\cref{prop:g10:reffunc:sqrt}、
\cref{prop:g10:reffunc:cube}）便足以用来刹车、平衡杠杆，
并推算行星的运行。压轴之作是开普勒
第三定律，可直接从太阳系的数据中读出。

\medskip
\noindent\textbf{第一部分——刹车的平方律。}
一条经验法则给出干燥路面上汽车的刹车距离：
$d = \left(\frac{v}{10}\right)^{2}$ 米，车速为 $v$ 千米/小时。
\begin{enumerate}
  \item 计算车速分别为 $30$、$50$、$90$ 和
        $130$ 千米/小时时的刹车距离。
  \item 车速加倍时，刹车距离乘以多少？
        试由平方\omterm{def:g10:functions:function}{函数}的伸缩性质加以解释，
        并以 $50 \to 100$ 千米/小时进行验证。
  \item 取证分析：刹车痕迹长 $50$ 米，公式认定司机的
        车速为多少（精确到千米/小时）？哪一个
        基本\omterm{def:g10:functions:function}{函数}给出了这一答案？在哪个\omterm{def:g10:functions:function}{定义域}上
        这一读数是唯一确定的
        （\cref{prop:g10:reffunc:sqrt}）？
  \item 比较车速每增加\emph{一}千米/小时所多出的距离，
        分别在低速与高速时：计算
        $d(31) - d(30)$ 与 $d(91) - d(90)$。
        两个结果体现了抛物线（\cref{prop:g10:reffunc:square}）的
        哪一条性质？
  \item 实际刹车还需计入反应时间：约一秒，
        这段时间内汽车行驶 $0.28v$ 米：
        $D(v) = 0.28v + \left(\frac{v}{10}\right)^2$。计算
        $D(50)$ 与 $D(130)$。在一次项与平方项中，
        哪一项在城市车速下占主导，
        哪一项在高速公路上占主导？
\end{enumerate}

\medskip
\noindent\textbf{第二部分——日常生活中的双曲线。}
\begin{enumerate}[resume]
  \item 杠杆平衡的条件是：两侧的力 $\times$ 距离
        在两侧相等。一个 $60$ 千克的孩子坐在离支点 $2$ 米处；
        另一侧距离为 $d$ 处的平衡力为
        $F(d) = \frac{120}{d}$（以千克当量计）。
        计算 $F(0.5)$、$F(1)$、$F(3)$。
  \item 反比例\omterm{def:g10:functions:function}{函数}（\cref{prop:g10:reffunc:inverse}）的变化情况
        对杠杆意味着什么？阿基米德的豪言
        “给我一个立足之地，我就能撬动地球”
        为何竟隐藏在双曲线的尾部？
        是什么禁止了 $d = 0$？
  \item 声音与光的强度随距离的\emph{平方}而衰减：
        距灯 $1$ 米处的强度为 $100$ 单位；求 $2$、$3$
        和 $10$ 米处的强度。用球面解释
        指数 $2$：光在距离 $d$ 处铺展于多大面积上
        （阿基米德墓碑问题，见初中分册）？
  \item 校外旅行租一辆大巴需 $600$ 欧元，
        由 $n$ 名参加者平均分摊。列表给出
        $n = 10, 20, 30$ 时的人均费用，求
        使人均费用不超过 $25$ 欧元的人数有多少，
        并说明随着 $n$ 增大，双曲线的形状承诺了什么、
        又排除了什么。
  \item 印制 $x$ 张海报需 $2x + 3$ 欧元（其中
        $3$ 欧元为固定准备费用）。证明每张海报的\emph{平均}费用为
        $a(x) = 2 + \frac3x$，描述其变化，并
        从经济角度解释其水平渐近线。
        （可对照\cref{exo:g10:reffunc:10}中的代数运算。）
\end{enumerate}

\medskip
\noindent\textbf{第三部分——开普勒的和谐。}
对每颗行星，设 $a$ 为其到太阳的平均距离（
以天文单位计，地球 $= 1$），$T$ 为其公转周期（以年计）。
数据：火星 $a = 1.52$，$T = 1.88$；木星 $a = 5.20$，
$T = 11.86$；土星 $a = 9.54$，$T = 29.4$。
\begin{enumerate}[resume]
  \item 计算地球、火星和木星的 $a^3$ 与 $T^2$。
        开普勒在 1618 年注意到了什么？
  \item 将该定律写成\omterm{def:g10:functions:function}{函数}形式：$T = \sqrt{a^3} = a\sqrt a$。
        用土星检验这一定律。
  \item 两个预测：一颗小行星绕日运行于 $a = 4$ 天文单位处——
        求其周期（结果为\omterm{def:g10:numbers:sets}{整数}）；一颗彗星每
        $27$ 年回归一次——求其平均距离
        （寻找一个完全立方数）。
  \item 该定律把哪三个基本\omterm{def:g10:functions:function}{函数}糅合在一起，
        以及其伸缩规律：当 $a$ 乘以 $4$ 时，
        $T$ 将如何变化？（以小行星与地球
        相比加以检验。）
  \item 开普勒从第谷·布拉赫的数据中读出这一定律，比牛顿
        以万有引力作出解释早了七十年。用一句话说明：
        这段史实对于在数表中识别基本\omterm{def:g10:functions:function}{函数}的力量，
        说明了什么？
\end{enumerate}

\medskip
\noindent\textbf{第四部分——龟兔赛跑。}
\begin{enumerate}[resume]
  \item 将 $\sqrt x$、$x$ 与 $x^2$ 分别在 $\intoo{0}{1}$
        与 $\intoo{1}{+\infty}$ 上排序
        （见\cref{prop:g10:reffunc:compare}），验证这两种
        排序在 $x = 0.25$ 与 $x = 4$ 处是否成立。
  \item 完整求解不等式 $\sqrt x > x$（对 $x \geq 0$，
        谨慎两边平方，并以\omterm{def:g10:algebra:signtable}{符号表}结束）。
  \item 引入立方\omterm{def:g10:functions:function}{函数}：将 $x$、$x^2$、$x^3$ 排序，
        分别在 $x = 0.5$ 与 $x = 2$ 处，并求出
        平方\omterm{def:g10:functions:function}{函数}与立方\omterm{def:g10:functions:function}{函数}相交的所有点。
  \item 两种储蓄方案在 $t$ 年后分别支付：方案 A 为 $100\sqrt t$
        欧元，方案 B 为 $10 t^2$ 欧元。证明 B 在 $t^3 = 100$
        时超过 A，并给出
        交叉时间（精确到月）。根式与幂的长期
        对比有何启示？
  \item 总结：对本题中的每一规律——刹车（$v^2$）、杠杆（$1/d$）、
        灯光（$1/d^2$）、开普勒（$a^{3/2}$）、
        方案 A（$\sqrt t$）——用一句话说明其\omterm{def:g10:functions:graph}{图像}的
        哪一特征承载了物理意义（陡升、
        铅直渐近线、扩散的球面、伸缩指数、
        平缓增长）。
        结论：基本\omterm{def:g10:functions:function}{函数}是字母表，
        自然界正是用它们书写的。
\end{enumerate}
\end{problem}
